\documentclass[11pt,a4paper]{article}

\usepackage[dutch]{babel}
\usepackage{amsmath,amssymb}
\usepackage[margin=2.2cm]{geometry}
\usepackage{enumitem}
\usepackage{tikz}
\usepackage{pgfplots}
\usepackage{needspace}

\pgfplotsset{compat=1.18}

% ======================
% Puntentelling
% ======================
\newcounter{totaalpunten}
\newcounter{verlpunten}

\newcommand{\punten}[1]{%
  \addtocounter{totaalpunten}{#1}%
  \makebox[0pt][l]{\hspace{-4em}\footnotesize #1p}%
}

\newcommand{\verlpunt}[1]{%
  \addtocounter{totaalpunten}{#1}%
  \addtocounter{verlpunten}{#1}%
  \makebox[0pt][l]{\hspace{-4em}\footnotesize #1p v}%
}

% ======================
% Opgave
% ======================
\newcommand{\opgave}[1]{%
\item \textbf{#1}\par
}

\setlist[enumerate,2]{label=\alph*)}

\begin{document}



\begin{enumerate}[label=\textbf{\arabic*.}, ref=\arabic*]
\item \opgave{Inhaalwerk Opgave 1}
Herleid.
\begin{enumerate}[label=\alph*), ref=\alph*]
\punten{2}
\item \( (a + b)^2 (3a + 2b) - 4(3a + 2b)(3a - 2b) \)
\verlpunt{2}
\item \( \left(-x^2 + x\right)^3 \)
\punten{2}
\item \( B = \frac{q^2 - q - 6}{q^2 - 4} \)
\verlpunt{2}
\item \( \frac{1}{5}c + \frac{b}{c} \cdot \frac{3c}{8} \)
\punten{2}
\item \( \frac{7}{y-4} - \frac{y+4}{y-2} \)
\punten{2}
\item \( \frac{5r^2}{r + 6} : \frac{r}{3r + 18} \)
\punten{2}
\item \( 4\sqrt{45} - \sqrt{2} \cdot \sqrt{32} \)
\punten{2}
\item \( (2\sqrt{7} - 3\sqrt{2})^2 \)
\punten{4}
\item \( \frac{6\sqrt{80}}{4\sqrt{5}} + \sqrt{\frac{25}{16}} + \sqrt{4 \frac{1}{4}} \)
\end{enumerate}
\end{enumerate}
\vspace{1cm}

% ======================


\end{document}

